\section{Theoretical Background}\label{sec:background}

\subsection{Lakehouse and Apache Iceberg}
A \emph{data warehouse} traditionally places data into closely governed tables for stable analytical serving: schemas are controlled, ETL (\emph{Extract -- Transform -- Load}) operations clean and shape data before it is loaded, and engines query the curated tables. A \emph{data lake}, in contrast, stores diverse files (CSV, JSON, Parquet, Avro, \dots) on inexpensive object storage and does not require a shared schema; ingest becomes easier but table-level governance can be lost. The \emph{lakehouse} pattern combines the two directions: data remains in open formats on a lake, while a table layer adds \emph{metadata}, versioned state, and transactional semantics that satisfy ACID properties (\emph{Atomicity, Consistency, Isolation, Durability}), so different engines can read a single consistent analytical result \parencite{armbrust2021lakehouse}.

\emph{Object storage} is a storage layer in which data is accessed by \emph{key} through an HTTP API -- typically compatible with the Amazon S3 interface -- rather than through traditional file paths; each object is a blob of bytes with attached metadata. In this artifact, MinIO supplies an S3-compatible endpoint for both checkpoints and Iceberg tables. Apache Iceberg is not a query processor but a \emph{table format} -- a specification that identifies which data files make up a table, which \emph{snapshot} (a table's state at a point in time) is current, and which \emph{manifest} files (intermediate metadata listing data files) describe table evolution \parencite{icebergdocs}. Iceberg distinguishes two delete object types: \emph{equality-delete}, which marks rows for deletion by key value, and \emph{position-delete}, which marks rows by position within a file. Both mechanisms support upsert-style updates. Consequently, when the report states that ``lakehouse output exists,'' the concrete evidence consists of Parquet data files and Iceberg metadata objects produced by the sink on MinIO; it does not assert that MinIO itself executes SQL.

\subsection{Stream processing, state, and materialized views}
In the pipeline of this study, an \emph{event} is one taxi-trip record replayed as a JSON message. Stream processing updates results as such events arrive rather than waiting for a complete batch. The abstraction that serves near-real-time results here is the \emph{materialized view} (MV): while an ordinary view is only a query definition re-evaluated at read time, a materialized view is a derived query result that the system maintains incrementally as inputs change; readers always observe the current aggregate without recomputing from raw records. For the taxi dataset, the MV is grouped by \texttt{borough} and \texttt{zone} and exposes the fields \texttt{trip\_count}, \texttt{total\_fare}, and \texttt{avg\_distance}.

The lookup join and grouped aggregation both require \emph{state} -- intermediate information that is retained across many events to keep running sums and counts per zone correct. A \emph{checkpoint} is a deliberately persisted snapshot of that state so an engine has a point from which processing can resume or recover after a disruption. Regarding semantics, a pipeline can be \emph{at-least-once} (each event is processed at least once, possibly duplicating it), \emph{at-most-once} (events may be lost), or \emph{exactly-once} (each event affects the result exactly once). A complementary property is \emph{idempotence} -- applying the same operation many times yields the same outcome -- which is often used at sinks to turn at-least-once delivery into effective exactly-once results. RisingWave is configured here to store checkpoint and state material in the MinIO \texttt{rw-checkpoint} bucket. That configuration explains where state is placed, but by itself it does not establish recovery behavior or end-to-end exactly-once delivery for this artifact.

\subsection{Kafka-compatible event streaming and RisingWave}
Kafka established the now-standard partitioned distributed-log model for streaming. In this model, a \emph{topic} is a named event stream (a ``pipe'' that holds messages of one kind), a \emph{partition} is a slice of a topic that provides local ordering and parallelism, and an \emph{offset} is the sequence position of a message within a partition \parencite{kreps2011kafka}. The \emph{producer} writes messages to a topic, while a \emph{consumer} reads them; consumers track the offsets they have processed so they can resume without re-reading. Continux uses Redpanda -- a C++ broker that is wire-compatible with the Kafka API -- as the broker that separates the Vector replay producer from the RisingWave consumer \parencite{redpandadocs}. This intermediate layer allows replay emission and streaming computation to progress at different rates: if RisingWave processes more slowly than Vector emits, messages remain buffered in the topic.

The topic is named \texttt{nyc-taxi-events} and is provisioned with three partitions and one replica. Three partitions characterize stream division and parallelism, while one replica means this run does not evaluate broker-node loss under a production high-availability arrangement -- if a broker node is lost, no in-topic copy survives. Because the collected evidence also does not reconcile offsets across the complete path, the report does not present exactly-once behavior as an experimentally verified end-to-end result.

RisingWave reads the Kafka-compatible event source, joins S3-compatible lookup data, maintains MVs, and writes to an Iceberg sink. RisingWave documentation identifies an Iceberg sink option for exactly-once delivery that is enabled by default while sink decoupling is enabled \parencite{risingwaveIceberg}. ``Sink decoupling'' decouples the sink's write schedule from the MV computation cadence to improve stability. This is a documented connector configuration capability, not an independently measured end-to-end result for this artifact.

\subsection{Blue/Green switching and dependencies}
\emph{Blue/Green deployment} is a deployment pattern in which two versions of the same service are maintained side by side: \emph{blue} is the version currently serving real traffic, while \emph{green} is a candidate replacement that is already prepared but is not publicly reachable. Once the candidate is judged ready, a \emph{cutover} switches the public access point from blue to green; when the change is purely a rename, that operation is called a \emph{swap}. Continux applies this pattern to MV definitions rather than to Kubernetes pods: \texttt{mv\_zone\_stats} is the stable public name that clients continue to call after the cutover; \texttt{mv\_zone\_stats\_blue} (baseline) retains all joined records, and \texttt{mv\_zone\_stats\_green} filters records with negative \texttt{fare\_amount} or \texttt{trip\_distance}.

RisingWave specifies \texttt{ALTER MATERIALIZED VIEW ... SWAP WITH ...} as a materialized-view name-swap operation \parencite{risingwaveSwap}. ``Atomic'' in this context concerns the public name operation as observed within the engine -- meaning the rename appears instantaneous to a client and exposes no intermediate state; it does not by itself establish that every dependent consumer or sink also changes its upstream definition. RisingWave's streaming-job modification guidance treats downstream dependencies such as sinks separately \parencite{risingwaveAlterJobs}; consequently, this study's cutover finding is bounded to the public query endpoint and does not extend to the downstream Iceberg sink.

\subsection{GitOps and observability}
OpenGitOps defines four principles for desired system state: it must be \emph{declarative} (describing ``what'' rather than ``how''), \emph{versioned} in Git, automatically \emph{pulled} by an agent, and continuously \emph{reconciled} -- the agent compares live resources with the desired configuration and drives them back toward that configuration \parencite{opengitops}. Continux uses Argo CD in the \emph{App-of-Apps} pattern, in which a root Application points to a directory containing child Applications; this layout lets the entire cluster be bootstrapped and synchronized from a single entry point \parencite{argocddocs}. Scaling Vector during a replay is a deliberate measurement action taken outside the GitOps path, not a replacement for Git as the long-lived configuration source.

\emph{Observability} is the ability to infer a system's internal condition from external signals such as metrics, logs, and workload state. A \emph{metric} is a numerical \emph{time series} -- a sequence of timestamped values representing a quantity over time; metrics are typically \emph{scraped} (pulled on a schedule) by a backend from an HTTP endpoint exposed either by the workload itself or by an \emph{exporter} (a side component that translates internal indicators into a standard metric format). VictoriaMetrics serves as the time-series backend in this artifact, storing metrics over time \parencite{victoriametricsdocs}; Grafana provides dashboard visualization \parencite{grafanadocs}; and the Continux exporter translates SQL counts and cutover results into standard metrics that VictoriaMetrics scrapes. Dashboards aid interpretation, but reported findings are always cross-checked against actual SQL output and MinIO object listings so that an incomplete or misleading panel is not treated as sole evidence.

\subsection{Ursa and related systems}
Ursa is a leaderless, Kafka-compatible, lakehouse-native streaming engine that targets cross-AZ replication, disk-storage, and connector cost drivers, while reporting throughput and semantics of its own design \parencite{merli2025ursa}. Continux neither implements Ursa nor compares cost against it; it adds an operational study of GitOps, dashboards, and public analytical-logic replacement in a small deployed stack.

Flink unifies stream and batch processing in one engine \parencite{carbone2015flink}; Kafka Streams and ksqlDB represent approaches aligned with Kafka ecosystems. RisingWave was chosen here for SQL/MV processing and the MV swap command. This positioning is conceptual and is not a comparative performance experiment.

\begin{table*}[t!]
\centering
\caption{Technology layers and evaluated roles}
\label{tab:tech-en}
\begin{tabularx}{\textwidth}{@{}p{0.20\textwidth}p{0.29\textwidth}Y@{}}
\toprule
\textbf{Layer} & \textbf{Technology} & \textbf{Evaluated role}\\
\midrule
Orchestration & K3s, Argo CD & Run workloads and reconcile manifests\\
Ingestion & Vector & Controlled JSONL replay\\
Event broker & Redpanda & Kafka-compatible event topic\\
Stream compute & RisingWave & Join, aggregate, public/blue/green MVs\\
Lakehouse storage & MinIO, Iceberg & Checkpoint and output table objects\\
Observation & VictoriaMetrics, Grafana, exporter & Counts, swap, health, and resources\\
\bottomrule
\end{tabularx}
\end{table*}
